\documentclass[11pt]{article}
\usepackage[utf8]{inputenc}
\usepackage[margin=1in]{geometry}
\usepackage{natbib}
\usepackage{booktabs}
\usepackage{amsmath}
\usepackage{graphicx}
\usepackage{hyperref}
\usepackage{xcolor}

\title{A SUMO E3 Ligase Promotes Long Non-Coding RNA Transcription for Small RNA-Directed DNA Elimination in \textit{Tetrahymena}}

\author{
  Author One$^{1}$, Author Two$^{1}$, Author Three$^{2}$, Author Four$^{1,*}$ \\
  $^{1}$Department of Molecular Biology, University A \\
  $^{2}$Department of Genetics, University B \\
  $^{*}$Corresponding author
}

\date{}

\begin{document}
\maketitle

\begin{abstract}
During sexual development in \textit{Tetrahymena thermophila}, Piwi-associated small RNAs (scnRNAs) guide programmed DNA elimination of germline-limited sequences from the developing somatic nucleus. This process requires nascent long non-coding RNAs (lncRNAs) transcribed from the parental somatic genome, which serve as templates for target-directed small RNA degradation (TDSD). However, the mechanism driving this lncRNA transcription has remained unclear. Here we show that the SUMO E3 ligase Ema2 is essential for somatic lncRNA accumulation, TDSD, and the production of viable sexual progeny. Ema2 knockout reduces lncRNA levels 6--10 fold and abolishes scnRNA degradation, with only $\sim$5\% progeny viability compared to $\sim$95\% in wild type. We demonstrate that Ema2 interacts with the SUMO E2 conjugating enzyme Ubc9 in a RING domain-dependent manner and promotes SUMOylation of the transcription elongation factor Spt6, reducing Spt6-SUMO levels to 15\% of wild type in Ema2-KO cells. Chromatin immunoprecipitation reveals that Ema2 is required for normal chromatin association of both Spt6 (reduced $\sim$4-fold) and RNA Polymerase II (reduced $\sim$3.3-fold) at lncRNA-producing loci. We propose a model in which Ema2-directed SUMOylation of Spt6 facilitates Pol~II loading onto chromatin to drive the lncRNA transcription that is a prerequisite for scnRNA-directed heterochromatin formation and DNA elimination.
\end{abstract}

\section{Introduction}

Programmed genome rearrangement is a widespread phenomenon in ciliated protozoa, where germline-limited DNA sequences are systematically eliminated from the developing somatic macronucleus during sexual reproduction \citep{chalker2011,betermier2014}. In \textit{Tetrahymena thermophila}, this process eliminates thousands of internally eliminated sequences (IESs) comprising approximately 30\% of the germline genome \citep{mochizuki2010}. The elimination is guided by a class of Piwi-associated small RNAs called scan RNAs (scnRNAs), which are produced from the germline micronuclear genome during meiosis \citep{mochizuki2002,aronica2008}.

The scnRNA pathway involves a crucial selection step called target-directed small RNA degradation (TDSD). During TDSD, scnRNAs are transported to the parental macronucleus, where they base-pair with nascent long non-coding RNAs (lncRNAs) transcribed from the somatic genome. ScnRNAs that find complementary sequences in somatic lncRNAs are degraded, while those without somatic matches---corresponding to germline-limited sequences---survive and are redirected to the developing macronucleus to guide DNA elimination \citep{noto2015,mochizuki2012}. Thus, active transcription of lncRNAs from the somatic genome is a prerequisite for scnRNA selection and subsequent DNA elimination.

Despite the central role of somatic lncRNA transcription in this pathway, the molecular machinery that actively promotes this transcription has remained poorly understood. Post-translational modification by the Small Ubiquitin-like Modifier (SUMO) is known to regulate transcription in diverse eukaryotes \citep{gareau2010,hendriks2017}. SUMOylation of transcription factors, chromatin remodelers, and RNA Polymerase II-associated factors has been shown to influence gene expression both positively and negatively \citep{rosonina2017,texari2018}. However, the involvement of SUMOylation in lncRNA-dependent genome rearrangement has not been explored.

Here we investigate the role of Ema2, a SUMO E3 ligase, in the scnRNA-directed DNA elimination pathway of \textit{Tetrahymena}. We demonstrate that Ema2 is essential for somatic lncRNA accumulation, TDSD, and the production of viable sexual progeny. Mechanistically, Ema2 interacts with the SUMO E2 conjugating enzyme Ubc9 and promotes SUMOylation of the transcription elongation factor Spt6, thereby facilitating the chromatin association of both Spt6 and RNA Polymerase II at lncRNA-producing loci.

\section{Related Work}

\subsection{Programmed DNA Elimination in Ciliates}

Programmed DNA elimination in ciliates has been studied most extensively in \textit{Tetrahymena} and \textit{Paramecium} \citep{chalker2011,duharcourt2009}. In both organisms, small RNA pathways guide the identification of sequences to be eliminated. The scnRNA pathway in \textit{Tetrahymena} was first described by \citet{mochizuki2002} and has since been shown to involve Argonaute/Piwi family proteins (Twi1), the RNA-dependent RNA polymerase Rdr1, and the Dicer-like protein Dcl1 \citep{mochizuki2010,aronica2008}. The TDSD mechanism was elucidated by \citet{noto2015}, who demonstrated that scnRNA selection occurs through interaction with somatic-origin lncRNAs.

\subsection{SUMOylation in Transcription Regulation}

SUMOylation is a reversible post-translational modification that regulates protein localization, stability, and activity \citep{gareau2010}. The SUMO pathway involves an E1 activating enzyme, an E2 conjugating enzyme (Ubc9), and E3 ligases that confer substrate specificity \citep{johnson2004}. In transcription, SUMOylation of chromatin-associated factors can promote or repress gene expression depending on the context \citep{rosonina2017}. Spt6, a histone chaperone and transcription elongation factor, has been identified as a SUMO target in several organisms, but the functional consequences of its SUMOylation have been unclear \citep{bortvin1996,close2011}.

\subsection{Long Non-Coding RNAs in Genome Defense}

Long non-coding RNAs serve as scaffolds and guides in diverse genome defense mechanisms, including piRNA-mediated transposon silencing in animals \citep{czech2018} and heterochromatin formation in fission yeast \citep{verdel2004}. In ciliates, lncRNAs provide the sequence information needed for scnRNA selection \citep{noto2015,mochizuki2012}. However, the transcriptional regulation of these lncRNAs has remained a gap in our understanding of ciliate genome rearrangement.

\section{Methods}

\subsection{Strains and Culture Conditions}

\textit{Tetrahymena thermophila} wild-type strains and Ema2 knockout strains were maintained and crossed according to standard protocols \citep{orias2011}. Ema2-KO strains were generated by replacing the Ema2 coding sequence with a drug resistance cassette via homologous recombination. Ema2-rescue strains were constructed by reintroducing the wild-type Ema2 gene at an ectopic locus.

\subsection{Progeny Viability Assays}

Mating pairs were isolated individually into drops of starvation medium and allowed to complete conjugation. Sexual progeny were identified by drug resistance markers and scored for viability by growth in selective medium. At least 100--150 pairs were scored per genotype.

\subsection{RT-qPCR for lncRNA Accumulation}

Total RNA was extracted at time points during conjugation (0, 3, 6, 9, and 12~h post-mixing). Somatic-origin lncRNAs were quantified by RT-qPCR using locus-specific primers, with normalization to rRNA levels.

\subsection{Co-Immunoprecipitation}

FLAG-tagged Ema2 (wild type or RING-domain mutant) and HA-tagged Ubc9 were co-expressed in \textit{Tetrahymena}. Cell lysates were immunoprecipitated with anti-FLAG antibodies and probed for HA-Ubc9 by western blot.

\subsection{In Vivo SUMOylation Assay}

Spt6 was immunoprecipitated from conjugating cells and probed with anti-SUMO antibodies. Band intensities were quantified by densitometry and normalized to the wild-type signal.

\subsection{Chromatin Immunoprecipitation (ChIP)}

ChIP-qPCR was performed at 8~h post-mixing using antibodies against Spt6 and RNA Pol~II (Rpb1 subunit). Enrichment was calculated as fold over IgG control at somatic lncRNA-producing loci.

\subsection{Small RNA Northern Blots}

Total small RNA was isolated from conjugating cells at 8~h and analyzed by northern blot to assess scnRNA levels as a readout of TDSD efficiency.

\section{Results}

\subsection{Ema2 Is Essential for Viable Sexual Progeny and DNA Elimination}

To determine whether Ema2 is required for sexual development, we crossed Ema2-KO strains and assessed progeny viability (Table~\ref{tab:viability}). Wild-type crosses yielded $\sim$95\% viable progeny (150 pairs), whereas Ema2-KO crosses produced only $\sim$5\% viable progeny, a near-complete loss of viability. Reintroduction of wild-type Ema2 (Ema2-rescue) restored viability to $\sim$90\% (100 pairs), confirming that the phenotype is attributable to Ema2 loss. DNA elimination was not completed in Ema2-KO progeny.

\begin{table}[ht]
\centering
\caption{Progeny viability in Ema2 knockout and rescue strains.}
\label{tab:viability}
\begin{tabular}{lrrr}
\toprule
\textbf{Genotype} & \textbf{Pairs Scored} & \textbf{Viable Progeny (\%)} & \textbf{DNA Elimination} \\
\midrule
WT $\times$ WT & 150 & $\sim$95 & Complete \\
Ema2-KO $\times$ Ema2-KO & 150 & $\sim$5 & Incomplete \\
Ema2-rescue $\times$ Ema2-rescue & 100 & $\sim$90 & Complete \\
\bottomrule
\end{tabular}
\end{table}

\subsection{Ema2 Is Required for Somatic lncRNA Accumulation}

RT-qPCR analysis revealed that somatic-origin lncRNAs accumulated dramatically during conjugation in wild-type cells, peaking at $\sim$18.7-fold above baseline at 9~h post-mixing (Table~\ref{tab:lncrna}). In contrast, Ema2-KO cells showed severely reduced lncRNA accumulation, reaching only $\sim$2.1-fold at the same time point---a 6--10 fold reduction compared to wild type. This demonstrates that Ema2 is required for the burst of somatic lncRNA transcription that normally occurs during conjugation.

\begin{table}[ht]
\centering
\caption{Somatic lncRNA accumulation during conjugation (fold change over baseline).}
\label{tab:lncrna}
\begin{tabular}{lrr}
\toprule
\textbf{Time (h)} & \textbf{WT} & \textbf{Ema2-KO} \\
\midrule
0 & 1.0 & 1.0 \\
3 & 4.5 & 1.2 \\
6 & 12.3 & 1.8 \\
9 & 18.7 & 2.1 \\
12 & 15.2 & 1.9 \\
\bottomrule
\end{tabular}
\end{table}

\subsection{TDSD Is Abolished in Ema2 Knockout Cells}

Since lncRNAs serve as substrates for TDSD, we assessed scnRNA degradation by northern blot at 8~h conjugation. In wild-type cells, $\sim$25\% of input scnRNAs remained at this time point, indicating efficient TDSD. In Ema2-KO cells, $\sim$85\% of scnRNAs remained, demonstrating that TDSD is severely impaired in the absence of Ema2-dependent lncRNA transcription.

\subsection{Ema2 Interacts with Ubc9 via Its RING Domain}

Co-immunoprecipitation experiments demonstrated that FLAG-Ema2 robustly co-precipitated HA-Ubc9, the sole SUMO E2 conjugating enzyme. A RING domain mutation (Ema2-RING-mut) abolished this interaction, confirming that the Ema2--Ubc9 interaction is RING domain-dependent and consistent with Ema2 functioning as a SUMO E3 ligase.

\subsection{Ema2 Promotes SUMOylation of Spt6}

In vivo SUMOylation assays revealed that Spt6 is SUMOylated in wild-type conjugating cells (relative band intensity 1.00). In Ema2-KO cells, Spt6 SUMOylation was reduced to 0.15 (an 85\% reduction), and Ema2-rescue restored SUMOylation to 0.88 (Table~\ref{tab:sumoylation}). These results establish Ema2 as the primary E3 ligase responsible for Spt6 SUMOylation during conjugation.

\begin{table}[ht]
\centering
\caption{Spt6 SUMOylation levels across genotypes.}
\label{tab:sumoylation}
\begin{tabular}{lr}
\toprule
\textbf{Condition} & \textbf{Spt6-SUMO Intensity (AU)} \\
\midrule
WT & 1.00 \\
Ema2-KO & 0.15 \\
Ema2-rescue & 0.88 \\
\bottomrule
\end{tabular}
\end{table}

\subsection{Ema2 Facilitates Spt6 and RNA Pol~II Chromatin Association}

ChIP-qPCR at somatic lncRNA-producing loci during conjugation (8~h) showed that both Spt6 and RNA Pol~II occupancy were substantially reduced in Ema2-KO cells (Table~\ref{tab:chip}). Spt6 enrichment was 8.5-fold over IgG in wild type but only 2.1-fold in Ema2-KO, a $\sim$4-fold reduction. RNA Pol~II (Rpb1) enrichment dropped from 11.3-fold to 3.4-fold, a $\sim$3.3-fold reduction. These results indicate that Ema2-dependent SUMOylation of Spt6 promotes the recruitment or stabilization of the transcription machinery at lncRNA-producing chromatin.

\begin{table}[ht]
\centering
\caption{ChIP-qPCR enrichment at somatic lncRNA loci (8~h conjugation).}
\label{tab:chip}
\begin{tabular}{lrr}
\toprule
\textbf{Factor} & \textbf{WT (fold/IgG)} & \textbf{Ema2-KO (fold/IgG)} \\
\midrule
Spt6 & 8.5 & 2.1 \\
RNA Pol~II (Rpb1) & 11.3 & 3.4 \\
\bottomrule
\end{tabular}
\end{table}

\section{Discussion}

Our findings establish a direct mechanistic link between SUMOylation and the lncRNA transcription that underlies small RNA-directed genome rearrangement in \textit{Tetrahymena}. We propose the following pathway: Ema2 (E3 ligase) partners with Ubc9 (E2) to SUMOylate the transcription elongation factor Spt6; SUMOylated Spt6 facilitates loading of both itself and RNA Pol~II onto chromatin at somatic lncRNA loci; the resulting lncRNA transcription provides the substrates for TDSD; and scnRNAs that survive TDSD guide heterochromatin formation and DNA elimination in the developing macronucleus.

The near-complete loss of progeny viability ($\sim$5\% vs. $\sim$95\%) in Ema2-KO crosses underscores the essential nature of this pathway. The dramatic reduction in lncRNA accumulation (6--10 fold) in Ema2-KO cells, combined with the failure of TDSD (85\% scnRNA retention vs. 25\%), provides a clear mechanistic explanation for this phenotype: without Ema2-driven lncRNA transcription, the scnRNA selection process cannot occur, and DNA elimination fails.

The identification of Spt6 as a key SUMOylation target of Ema2 is particularly informative. Spt6 is a histone chaperone that facilitates transcription elongation by maintaining chromatin structure \citep{bortvin1996,close2011}. SUMOylation has been shown to regulate the chromatin association of various transcription factors in other systems \citep{rosonina2017}, and our ChIP data indicate that Spt6 SUMOylation by Ema2 serves a similar function at lncRNA loci. The concomitant reduction in Pol~II occupancy suggests that Spt6 may act upstream of Pol~II recruitment, consistent with models in which histone chaperones prepare chromatin for transcription \citep{ivanovska2011}.

These findings add SUMOylation to the growing list of post-translational modifications involved in ciliate genome rearrangement, which already includes histone methylation (H3K9me3 and H3K27me3) \citep{liu2007,taverna2002}. Whether Ema2 has additional substrates beyond Spt6 remains to be determined. Future proteomic studies may reveal a broader SUMOylation program that coordinates the multiple chromatin-level events required for DNA elimination.

\section{Conclusion}

We have identified the SUMO E3 ligase Ema2 as an essential regulator of somatic lncRNA transcription in \textit{Tetrahymena}. Ema2 promotes SUMOylation of the transcription elongation factor Spt6, facilitating the chromatin association of Spt6 and RNA Pol~II at lncRNA-producing loci. This mechanism ensures the production of lncRNAs required for scnRNA selection via TDSD, which is in turn required for programmed DNA elimination and viable sexual progeny. Our work reveals SUMOylation as a critical regulatory layer in small RNA-directed genome rearrangement.

\bibliographystyle{plainnat}
\bibliography{references}

\end{document}
